% Options for packages loaded elsewhere
\PassOptionsToPackage{unicode}{hyperref}
\PassOptionsToPackage{hyphens}{url}
\documentclass[
  12pt,
]{ctexart}
\usepackage{xcolor}
\usepackage{amsmath,amssymb}
\setcounter{secnumdepth}{-\maxdimen} % remove section numbering
\usepackage{iftex}
\ifPDFTeX
  \usepackage[T1]{fontenc}
  \usepackage[utf8]{inputenc}
  \usepackage{textcomp} % provide euro and other symbols
\else % if luatex or xetex
  \usepackage{unicode-math} % this also loads fontspec
  \defaultfontfeatures{Scale=MatchLowercase}
  \defaultfontfeatures[\rmfamily]{Ligatures=TeX,Scale=1}
\fi
\usepackage{lmodern}
\ifPDFTeX\else
  % xetex/luatex font selection
\fi
% Use upquote if available, for straight quotes in verbatim environments
\IfFileExists{upquote.sty}{\usepackage{upquote}}{}
\IfFileExists{microtype.sty}{% use microtype if available
  \usepackage[]{microtype}
  \UseMicrotypeSet[protrusion]{basicmath} % disable protrusion for tt fonts
}{}
\makeatletter
\@ifundefined{KOMAClassName}{% if non-KOMA class
  \IfFileExists{parskip.sty}{%
    \usepackage{parskip}
  }{% else
    \setlength{\parindent}{0pt}
    \setlength{\parskip}{6pt plus 2pt minus 1pt}}
}{% if KOMA class
  \KOMAoptions{parskip=half}}
\makeatother
\usepackage{longtable,booktabs,array}
\usepackage{calc} % for calculating minipage widths
% Correct order of tables after \paragraph or \subparagraph
\usepackage{etoolbox}
\makeatletter
\patchcmd\longtable{\par}{\if@noskipsec\mbox{}\fi\par}{}{}
\makeatother
% Allow footnotes in longtable head/foot
\IfFileExists{footnotehyper.sty}{\usepackage{footnotehyper}}{\usepackage{footnote}}
\makesavenoteenv{longtable}
\setlength{\emergencystretch}{3em} % prevent overfull lines
\providecommand{\tightlist}{%
  \setlength{\itemsep}{0pt}\setlength{\parskip}{0pt}}
\usepackage{bookmark}
\IfFileExists{xurl.sty}{\usepackage{xurl}}{} % add URL line breaks if available
\urlstyle{same}
\hypersetup{
  hidelinks,
  pdfcreator={LaTeX via pandoc}}

\author{SCX}
\date{2026}

\begin{document}

\section{SCX 论文框架 TODO}\label{scx-ux8bbaux6587ux6846ux67b6-todo}

\begin{quote}
创建日期：2026-06-26 基于：五审查员综合报告
关联文件：{[}{[}论文框架审查报告\_五审查员综合{]}{]}
\end{quote}

\begin{center}\rule{0.5\linewidth}{0.5pt}\end{center}

\subsection{阶段
0：立即止损（本周）}\label{ux9636ux6bb5-0ux7acbux5373ux6b62ux635fux672cux5468}

\subsubsection{0.1 EGP Paper 1
内容收窄}\label{egp-paper-1-ux5185ux5bb9ux6536ux7a84}

\begin{itemize}
\tightlist
\item[$\square$]
  \textbf{移除 EGP Paper 1 中的 SCX 噪声检测 Figs (Fig 4-8)}

  \begin{itemize}
  \tightlist
  \item
    这些图（SCX 噪声分布、一层 vs 两层 F1、fmax vs error、综合力
    RMSE、雷达图）属于 SCX-MLIP
  \item
    EGP P1 只保留：EOS 对比、弹性常数对比、Per-batch 误差对比、Gauge
    violation 对比
  \item
    参考：\texttt{G:\textbackslash{}Xiaogan\_Supercomputing\_data\textbackslash{}SCX\textbackslash{}paper\textbackslash{}paper1\_mlip\textbackslash{}FIGURES\_ANALYSIS.md}
  \end{itemize}
\item[$\square$]
  \textbf{移除 EGP P1 中所有 SCX 理论提及}

  \begin{itemize}
  \tightlist
  \item
    不在正文中讨论''状态条件专家可靠性''
  \item
    Discussion 末尾最多一句话：``Whether expert reliability is
    inherently state-conditioned remains an open question for future
    work''
  \end{itemize}
\item[$\square$]
  \textbf{确认 EGP P1 的独立贡献叙事}

  \begin{itemize}
  \tightlist
  \item
    只 claim：gauge-fixed shared+correction ACE 参数化
  \item
    不 claim：新的 MoE runtime、element correction 原子能、AlGaN 可转移
  \end{itemize}
\end{itemize}

\subsubsection{0.2 SCX 理论加固}\label{scx-ux7406ux8bbaux52a0ux56fa}

\begin{itemize}
\tightlist
\item[$\square$]
  \textbf{删除 Arrow
  不可能定理类比}（\texttt{theory/propositions/01\_global\_ranking\_insufficiency.md}
  Section 4）
\item[$\square$]
  \textbf{重构 Proposition 1}

  \begin{itemize}
  \tightlist
  \item
    从''不存在全局最优专家''改为''全局排序的遗憾下界为 Ω(p·δ\_min)''
  \item
    给出可检验的共单调性条件（已有推论草稿，需正式化）
  \end{itemize}
\item[$\square$]
  \textbf{添加不可识别性定理草稿}

  \begin{itemize}
  \tightlist
  \item
    显式陈述：无额外假设时，噪声与可学习困难不能仅从 (x,y) 识别
  \item
    放在 Prop 2 旁边，作为理论边界
  \end{itemize}
\end{itemize}

\begin{center}\rule{0.5\linewidth}{0.5pt}\end{center}

\subsection{阶段 1：EGP Paper 1 收尾（1-2
周）}\label{ux9636ux6bb5-1egp-paper-1-ux6536ux5c3e1-2-ux5468}

\subsubsection{1.1 写作}\label{ux5199ux4f5c}

\begin{itemize}
\tightlist
\item[$\square$]
  \textbf{写完 Introduction}

  \begin{itemize}
  \tightlist
  \item
    锚定问题：多元素 ACE 势函数合并时 gauge 不一致导致势能面不连续
  \item
    文献回顾：ACE 框架 (Drautz 2019, Lysogorskiy 2021)，PACE，已有
    merging 尝试
  \item
    贡献陈述：shared+correction 参数化 + coefficient-level post-hoc
    gauge fixing
  \end{itemize}
\item[$\square$]
  \textbf{写完 Methods}

  \begin{itemize}
  \tightlist
  \item
    2.1: Shared+Correction ACE Architecture (E = Σ{[}c₀ᵀB + c\_ZᵀB +
    b\_Z{]})
  \item
    2.2: Gauge Freedom (c₀→c₀+d, c\_Z→c\_Z-d 不变性)
  \item
    2.3: Post-hoc Projection (g = Σπ\_Z c\_Z; c\_Z' = c\_Z - g; c₀' = c₀
    + g)
  \item
    2.4: Why Soft Constraint Fails (梯度竞争分析：loss 范比
    79\%→5.6\%，精度崩溃)
  \item
    2.5: Training Protocol (AlN v3: Pacemaker baseline → Model A
    ablation → Model B design matrix → constrained ridge → export PACE)
  \end{itemize}
\item[$\square$]
  \textbf{写完 Results: AlN v3}

  \begin{itemize}
  \tightlist
  \item
    3.1: EOS (V₀, B₀ vs DFT ref 10.64/192.6)
  \item
    3.2: Elastic Constants (C₁₁-C₆₆ 对比表)
  \item
    3.3: Phonon Forces (声子力 RMSE 0.00793→0.00580, -27\%)
  \item
    3.4: Gauge Violation (8.77→4.6×10⁻¹⁶, soft constraint 失败 vs
    post-hoc 成功)
  \item
    3.5: Surface/Defect Transferability
  \end{itemize}
\item[$\square$]
  \textbf{写完 Discussion + Conclusion}
\item[$\square$]
  \textbf{制作最终 Figures}（EOS curve, elastic bar chart, gauge
  violation comparison）
\end{itemize}

\subsubsection{1.2 实验补充（可选，取决于 GaN/AlN v4
进度）}\label{ux5b9eux9a8cux8865ux5145ux53efux9009ux53d6ux51b3ux4e8e-ganaln-v4-ux8fdbux5ea6}

\begin{itemize}
\tightlist
\item[$\square$]
  \textbf{检查 GaN/AlN v4 超算任务状态}

  \begin{itemize}
  \tightlist
  \item
    \texttt{python\ .hpc/hpcmgr.py\ list}
  \item
    \texttt{python\ .hpc/hpcmgr.py\ logs\ \textless{}job\_id\textgreater{}}
  \end{itemize}
\item[$\square$]
  \textbf{如果 v4 已完成}：补充 GaN/AlN v4 的 Model B 训练 + gauge fix +
  对比
\item[$\square$]
  \textbf{如果 v4 未完成}：EGP P1 先以 AlN v3 单系统投出，v4 结果作为
  revision 补充
\end{itemize}

\subsubsection{1.3 投稿准备}\label{ux6295ux7a3fux51c6ux5907}

\begin{itemize}
\tightlist
\item[$\square$]
  \textbf{选择目标期刊}：npj Computational Materials（首选）/
  PRM（备选）/ JCTC（备选）
\item[$\square$]
  \textbf{准备 SI（Supporting Information）}
\item[$\square$]
  \textbf{准备 Cover Letter}
\item[$\square$]
  \textbf{arXiv 上传}
\end{itemize}

\begin{center}\rule{0.5\linewidth}{0.5pt}\end{center}

\subsection{阶段 2：P0 实验 ---
重训验证（本周，不可跳过）⭐}\label{ux9636ux6bb5-2p0-ux5b9eux9a8c-ux91cdux8badux9a8cux8bc1ux672cux5468ux4e0dux53efux8df3ux8fc7}

\subsubsection{2.1 AlN v3
去噪重训}\label{aln-v3-ux53bbux566aux91cdux8bad}

\begin{itemize}
\tightlist
\item[$\square$]
  \textbf{准备去噪数据集}

  \begin{itemize}
  \tightlist
  \item
    方案 A：移除 fmax \textgreater{} 5 eV/Å 的 74 帧（保守）
  \item
    方案 B：移除 SCX 两层 Top-2 噪声状态（约 90-131 帧）
  \item
    方案 C：移除 fmax \textgreater{} 10 eV/Å 的 14 帧 + 降权 fmax 5-10
    eV/Å 的 60 帧
  \item
    推荐先做方案 A（最简单，审稿人最容易理解）
  \end{itemize}
\item[$\square$]
  \textbf{训练对照组}

  \begin{itemize}
  \tightlist
  \item
    \begin{enumerate}
    \def\labelenumi{(\alph{enumi})}
    \tightlist
    \item
      Clean baseline: 去噪后数据训练 Single ACE
    \end{enumerate}
  \item
    \begin{enumerate}
    \def\labelenumi{(\alph{enumi})}
    \setcounter{enumi}{1}
    \tightlist
    \item
      Original baseline: 全部 534 帧训练 Single ACE（已有）
    \end{enumerate}
  \item
    \begin{enumerate}
    \def\labelenumi{(\alph{enumi})}
    \setcounter{enumi}{2}
    \tightlist
    \item
      Random drop: 随机去除 74 帧，训练 Single ACE
    \end{enumerate}
  \item
    \begin{enumerate}
    \def\labelenumi{(\alph{enumi})}
    \setcounter{enumi}{3}
    \tightlist
    \item
      Loss-based drop: 去除训练 loss 最高的 74 帧，训练 Single ACE
    \end{enumerate}
  \end{itemize}
\item[$\square$]
  \textbf{评测}

  \begin{itemize}
  \tightlist
  \item
    在同一个 103 帧测试集上测力 RMSE、能量 RMSE
  \item
    报告 (a) vs (b) 的改善幅度（实测值，不是预估值！）
  \item
    展示 SCX 指导的 drop (a) 显著优于 random drop (c) 和 loss-based drop
    (d)
  \end{itemize}
\item[$\square$]
  \textbf{更新 FIGURES\_ANALYSIS.md 的 Fig 7}

  \begin{itemize}
  \tightlist
  \item
    用实测值替换''SCX-ACE（预估）``柱
  \end{itemize}
\end{itemize}

\subsubsection{2.2
重训后叙事更新}\label{ux91cdux8badux540eux53d9ux4e8bux66f4ux65b0}

\begin{itemize}
\tightlist
\item[$\square$]
  \textbf{根据实测结果确定 claim 级别}

  \begin{itemize}
  \tightlist
  \item
    如果力 RMSE 降 \textgreater15\%：可以 claim ``SCX improves model
    performance'' (L2)
  \item
    如果力 RMSE 降 5-15\%：claim ``SCX identifies data quality issues
    that affect model performance'' (L1+)
  \item
    如果力 RMSE 降 \textless5\%：claim ``SCX is a diagnostic tool for
    data quality auditing'' (L1)
  \item
    如果力 RMSE 没降甚至升了：需要重新审视 SCX 的噪声检测逻辑
  \end{itemize}
\end{itemize}

\begin{center}\rule{0.5\linewidth}{0.5pt}\end{center}

\subsection{阶段 3：SCX-MLIP 论文准备（2-4
周）}\label{ux9636ux6bb5-3scx-mlip-ux8bbaux6587ux51c6ux59072-4-ux5468}

\subsubsection{3.1 理论部分（精简到 \textasciitilde3
页正文）}\label{ux7406ux8bbaux90e8ux5206ux7cbeux7b80ux5230-3-ux9875ux6b63ux6587}

\begin{itemize}
\tightlist
\item[$\square$]
  \textbf{写 Core Definitions 节}

  \begin{itemize}
  \tightlist
  \item
    R\_m(s) = E{[}ℓ(f\_m(x), f*(x)) \textbar{} x∈s{]}
  \item
    SCX\_m(s) = P(ℓ \textless{} τ \textbar{} x∈s)
  \item
    数据四分类：Valuable / Noisy / Redundant / Expert-dependent
  \item
    不包含 V(s) 乘法分解（用四分类规则替代）
  \end{itemize}
\item[$\square$]
  \textbf{写 Three Propositions 节（重构版）}

  \begin{itemize}
  \tightlist
  \item
    Prop 1：全局排序的遗憾下界（含共单调性可检验条件）
  \item
    Prop 2：噪声下 error-driven sampling 次优性 + 不可识别性定理
  \item
    Prop 3：状态条件权重优势（Jensen + 塔性质证明）
  \end{itemize}
\item[$\square$]
  \textbf{写 State Discovery 节}

  \begin{itemize}
  \tightlist
  \item
    两层架构：L1 域知识 + L2 ErrorDrivenEncoder
  \item
    最优粒度：Silhouette + 风险同质性联合优化
  \item
    状态发现作为 oracle（ErrorDrivenEncoder 是实例化，细节放附录）
  \end{itemize}
\item[$\square$]
  \textbf{写冷启动协议}（可放附录）

  \begin{itemize}
  \tightlist
  \item
    第一阶段：无标签代理（expert disagreement, model uncertainty,
    self-consistency）
  \item
    第二阶段：acquisition
  \item
    收敛条件
  \end{itemize}
\end{itemize}

\subsubsection{3.2 实验部分}\label{ux5b9eux9a8cux90e8ux5206}

\begin{itemize}
\tightlist
\item[$\square$]
  \textbf{AlN v3 两层状态发现}（已有，整理成论文格式）

  \begin{itemize}
  \tightlist
  \item
    Table: 一层 vs 两层 F1 (th=2.0/3.0/4.0/5.0)
  \item
    Table: Top-K 噪声捕获率
  \item
    Figure: 50\% megastate 被分解为 6 个误差特征不同的子状态
  \item
    Figure: Phonon 100\% 落入同一低误差状态
  \end{itemize}
\item[$\square$]
  \textbf{AlN v3 去噪重训}（P0 完成后纳入）

  \begin{itemize}
  \tightlist
  \item
    Figure: 力 RMSE 实测对比（Clean vs Original vs Random vs
    Loss-based）
  \item
    Figure: fmax vs test error scatter (r=0.966)
  \end{itemize}
\item[$\square$]
  \textbf{跨材料验证}（P1）

  \begin{itemize}
  \tightlist
  \item
    下载 Si、Cu、MgO 公开 ACE 训练数据
  \item
    运行 SCX 两层分析
  \item
    报告每个系统的噪声/冗余发现
  \item
    目标：展示 SCX 不只在 AlN 上有效
  \end{itemize}
\item[$\square$]
  \textbf{合成实验}（P1）

  \begin{itemize}
  \tightlist
  \item
    已知噪声结构 → SCX 恢复状态结构的能力
  \item
    不可识别性定理的数值演示
  \end{itemize}
\end{itemize}

\subsubsection{3.3 竞争对比}\label{ux7adeux4e89ux5bf9ux6bd4}

\begin{itemize}
\tightlist
\item[$\square$]
  \textbf{与 loss-based sampling 对比}（包含在重训实验中）
\item[$\square$]
  \textbf{与 random sampling 对比}（包含在重训实验中）
\item[$\square$]
  \textbf{（可选）与 active learning 方法对比}（如 BADGE, uncertainty
  sampling）

  \begin{itemize}
  \tightlist
  \item
    仅当重训结果非常好且想冲高影响力期刊时做
  \end{itemize}
\end{itemize}

\subsubsection{3.4 写作与投稿}\label{ux5199ux4f5cux4e0eux6295ux7a3f}

\begin{itemize}
\tightlist
\item[$\square$]
  \textbf{写 Introduction}

  \begin{itemize}
  \tightlist
  \item
    锚定：MLIP 数据贵但价值不均 → EGP P1 的 9.7× per-batch 差异 →
    需要一个状态条件框架
  \item
    不泛化到''通用框架''------收窄到 ``methodology for MLIP data
    quality''
  \end{itemize}
\item[$\square$]
  \textbf{写 Methods: SCX Pipeline for MLIP}
\item[$\square$]
  \textbf{写 Discussion}

  \begin{itemize}
  \tightlist
  \item
    SCX 是数据诊断工具，与模型架构改进互补
  \item
    拓展潜力（一句话提 Sim/Health，不展开）
  \item
    局限：冷启动需要初始标签; 状态粒度敏感
  \end{itemize}
\item[$\square$]
  \textbf{内部审查 → 修改 → arXiv → 投稿}
\end{itemize}

\begin{center}\rule{0.5\linewidth}{0.5pt}\end{center}

\subsection{阶段 4：跨领域扩展（SCX-MLIP
投稿后）}\label{ux9636ux6bb5-4ux8de8ux9886ux57dfux6269ux5c55scx-mlip-ux6295ux7a3fux540e}

\subsubsection{4.1
选择主攻领域}\label{ux9009ux62e9ux4e3bux653bux9886ux57df}

\begin{itemize}
\tightlist
\item[$\square$]
  \textbf{评估 SCX-Sim 可行性}

  \begin{itemize}
  \tightlist
  \item
    是否有产业合作方数据？
  \item
    是否能用公开仿真 benchmark？
  \item
    如果 YES → 优先 SCX-Sim（Nature Computational Science 目标）
  \end{itemize}
\item[$\square$]
  \textbf{评估 SCX-Health 可行性}

  \begin{itemize}
  \tightlist
  \item
    MedMNIST 实验 GPU 升级后是否有效？
  \item
    是否有临床合作方？
  \item
    如果 Sim 不可行 → 做 Health
  \end{itemize}
\end{itemize}

\subsubsection{4.2 领域深挖}\label{ux9886ux57dfux6df1ux6316}

\begin{itemize}
\tightlist
\item[$\square$]
  \textbf{选定的领域内，做出与 AlN v3 可比证据强度的实验结果}
\item[$\square$]
  \textbf{写第三篇论文}
\end{itemize}

\begin{center}\rule{0.5\linewidth}{0.5pt}\end{center}

\subsection{阶段
5：长期（论文系列完成后）}\label{ux9636ux6bb5-5ux957fux671fux8bbaux6587ux7cfbux5217ux5b8cux6210ux540e}

\subsubsection{5.1 理论完善}\label{ux7406ux8bbaux5b8cux5584}

\begin{itemize}
\tightlist
\item[$\square$]
  将 SCX 形式化为完整的统计学习理论论文（如果 SCX-MLIP 反响好）
\item[$\square$]
  V(s) 从信息增益第一性原理推导（如果找到合适的数学框架）
\end{itemize}

\subsubsection{5.2 开源}\label{ux5f00ux6e90}

\begin{itemize}
\tightlist
\item[$\square$]
  开源 \texttt{scx-research/} 代码（``足以复现论文''级别）
\item[$\square$]
  发布 2-3 个 SCX-cleaned 公开数据集（MPTraj 子集、ColabFit 的 AlN/GaN
  等）
\end{itemize}

\subsubsection{5.3 商业}\label{ux5546ux4e1a}

\begin{itemize}
\tightlist
\item[$\square$]
  基于 SCX-Pro 的商业授权/合作谈判
\item[$\square$]
  SCX Consortium 产业联盟（远期）
\end{itemize}

\begin{center}\rule{0.5\linewidth}{0.5pt}\end{center}

\subsection{优先级总览}\label{ux4f18ux5148ux7ea7ux603bux89c8}

\begin{verbatim}
本周必做 ████████████████████ P0
████ 阶段 0: 止损（EGP P1 收窄 + 理论加固）
████ 阶段 2: AlN v3 去噪重训（4-8h 计算）

两周内 ████████████████████ P1
████ 阶段 1: EGP P1 写完 + 投稿
████ 阶段 2: 公开 MLIP 数据集分析

一月内 ████████████████████ P2
████ 阶段 3: SCX-MLIP 论文完成

投稿后 ████████████████████ P3
████ 阶段 4: 跨领域扩展（Sim 或 Health 二选一）
\end{verbatim}

\begin{center}\rule{0.5\linewidth}{0.5pt}\end{center}

\subsection{关键决策记录}\label{ux5173ux952eux51b3ux7b56ux8bb0ux5f55}

\begin{longtable}[]{@{}
  >{\raggedright\arraybackslash}p{(\linewidth - 4\tabcolsep) * \real{0.3333}}
  >{\raggedright\arraybackslash}p{(\linewidth - 4\tabcolsep) * \real{0.3333}}
  >{\raggedright\arraybackslash}p{(\linewidth - 4\tabcolsep) * \real{0.3333}}@{}}
\toprule\noalign{}
\begin{minipage}[b]{\linewidth}\raggedright
日期
\end{minipage} & \begin{minipage}[b]{\linewidth}\raggedright
决策
\end{minipage} & \begin{minipage}[b]{\linewidth}\raggedright
理由
\end{minipage} \\
\midrule\noalign{}
\endhead
\bottomrule\noalign{}
\endlastfoot
2026-06-26 & SCX-Theory 不与 SCX-MLIP 分开投稿 &
五审查员全票：独立理论论文被攻击风险极高，且 EGP P1
的保护不传导到不同审稿人 \\
2026-06-26 & 理论部分收窄到 1/3 篇幅 &
核心期刊审稿人更看重实验验证，长理论会稀释实验贡献的冲击力 \\
2026-06-26 & EGP P1 不包含 SCX 噪声检测 &
避免审稿人混淆两篇论文的贡献边界 \\
2026-06-26 & 七领域收窄到 MLIP + 最多一个 & 当前 MedMNIST/CIFAR
实验质量不足以支撑通用性 claim \\
\end{longtable}

\begin{center}\rule{0.5\linewidth}{0.5pt}\end{center}

\emph{关联文件：{[}{[}论文框架审查报告\_五审查员综合{]}{]}}

\end{document}
